\chapter{Price the Outcome}

An adviser in Miami helps companies navigate regulatory and clinical pathways.
When he explains his fee, he begins with the calendar. Suppose, he says, a
product is expected to produce \$12 million in annual sales. If his work brings
the launch forward by three months, the company receives three months of sales
it would otherwise have missed. Against that claimed \$3 million, a \$300,000
fee can appear small.

This is more useful than adding up the adviser's hours and applying a markup.
It asks what changes for the buyer. It is also only the beginning of a sound
calculation. Three months of a \$12 million annual run rate is \$3 million of
\emph{revenue}, not necessarily \$3 million of value. The launch may ramp
slowly. Some purchases may simply move forward in time. Production, selling,
and support consume cash. The product may fail to gain approval or demand. A
faster pathway has value only if the pathway remains lawful, clinically sound,
and safe.

Price should therefore begin with the outcome and then survive the details.
The buyer needs a credible account of the world with the offer, the world
without it, and the evidence that connects the two. The seller needs a price
that can support competent delivery. Neither side is served by impressive
arithmetic built on an imaginary result.

\section{Name the change being bought}

A product is an object or a service from the seller's side. From the buyer's
side, it is a proposed change. The buyer may be purchasing fewer errors, an
earlier opening, more usable capacity, a shorter queue, a lower probability of
loss, a more trusted experience, or the option to do something that was
previously impossible.

The distinction matters because effort and outcome do not move together. Ten
hours spent on the right bottleneck may be more valuable than a year spent on
work the buyer does not need. A beautifully made feature can have no commercial
value for this customer. A modest integration can become essential if it
removes the final obstacle to a large decision.

Begin with a sentence that contains a person, a baseline, a change, and a
period:

\begin{quote}
For this buyer, compared with the present alternative, the offer is expected
to produce this measurable change over this period.
\end{quote}

Each phrase does work. ``This buyer'' prevents a general benefit from being
assigned to everyone. The baseline prevents the seller from comparing the
offer with nothing. ``Expected'' leaves room for uncertainty. The period keeps
a temporary improvement from being treated as a perpetual gain.

Return to the Miami example. The real baseline is not simply a launch three
months later. It includes the buyer's lawful internal route, other qualified
advisers, the probability of approval under each route, the operational work
required after approval, and the value of being early in that particular
market. If the adviser changes none of those things, the calendar claim has
little foundation. If credible evidence shows that he can shorten the path
without weakening it, time becomes a legitimate part of the offer.

\section{Estimate value without disguising uncertainty}

A useful model separates the kinds of change before combining them. For a
particular buyer and use, write

\begin{equation}
\begin{aligned}
V_{\mathrm{buyer}}={}&
  \mathbb{E}[\text{incremental gain}]
  + \mathbb{E}[\text{loss avoided}]\\
 &+ V_{\mathrm{time}} + V_{\mathrm{option}} \\
 &- \text{implementation burden}\\
 &- \mathbb{E}[\text{downside caused}].
\end{aligned}
\end{equation}

This is a checklist expressed as an equation, not a machine for producing an
authoritative number. Each term needs a unit, a period, a baseline, and
evidence. Benefits that overlap should not be counted twice. A saved hour
cannot be valued as both reduced payroll and new productive capacity unless
the buyer can actually realize both.

For an accelerated launch, the time term might be estimated more carefully as

\begin{equation}
V_{\mathrm{time}}
  \approx p\,\frac{\Delta t}{12}
  \bigl(R m - C_{\mathrm{added}}\bigr)
  - C_{\mathrm{switch}},
\end{equation}

where \(p\) is the probability that the acceleration succeeds, \(\Delta t\) is
the number of months genuinely gained, \(R\) is annual incremental revenue,
\(m\) is the contribution margin on that revenue, and the final terms capture
added operating and switching costs. A ramp-up curve, cannibalized sales, or a
short-lived advantage may reduce the result further.

The calculation will often be a range rather than a point. That is a virtue.
Use a conservative case, a central case, and an upside case, then show which
assumption changes each one. If the offer is attractive only when every input
takes its most favorable value, the price rests on hope. If it remains useful
under a plausible conservative case, both parties have something firmer to
discuss.

Price also has two different boundaries. The seller's floor must support
delivery: direct cost, scarce capacity, expected rework or liability, and a
margin that keeps the promise serviceable. The buyer's ceiling is constrained
by risk-adjusted value, available alternatives, implementation burden, budget,
and the value of waiting. A possible agreement lies where those ranges overlap:

\begin{equation}
\begin{aligned}
P_{\mathrm{floor}}={}&C_{\mathrm{delivery}}+C_{\mathrm{capacity}}\\
&+C_{\mathrm{risk}}+M_{\mathrm{sustainable}},\\
P_{\mathrm{ceiling}}={}&V_{\mathrm{buyer}}
  -V_{\mathrm{best\ alternative}},\\
P_{\mathrm{floor}}\ &\leq P\leq P_{\mathrm{ceiling}}.
\end{aligned}
\end{equation}

The right side is deliberately simplified. A budget can prevent a sensible
purchase. Bargaining power can move the price. Procurement rules, timing,
financing, taxes, and switching costs can change the feasible range. The model
does not announce what is fair; it reveals what the proposed number assumes.

\section{The number belongs to a buyer}

One Silicon Valley founder gives an extreme example. He reports selling a
point-of-sale company for a little more than \$70 million even though it had
made only about \$60,000 in revenue and consumed about \$17 million. His
explanation is fit. The acquirer, he says, had 11,000 sales representatives and
needed capabilities his company had built. He estimated that hiring another
9,000 representatives represented roughly \$70--80 million to the acquirer,
and the parties settled on a number in that neighborhood.

The account is incomplete and self-reported. We do not know the transaction
terms, liabilities, intellectual property, competitive alternatives, retention
agreements, tax treatment, or whether the stated hiring comparison was the
decisive calculation. It should not teach founders that revenue is irrelevant
or that money already spent determines value. The \$17 million was sunk; the
low revenue also signaled uncertainty that a buyer could reasonably care about.

What the story does show is that value can be specific to a combination. A
technology may save one acquirer years of integration, fit its sales system,
block a competitor, or open an option unavailable to anyone else. Another
buyer may have no use for it. A strategic acquisition is therefore a poor
benchmark for the ordinary customer and a powerful reminder to identify the
actual buyer.

The same principle applies at smaller scale. A scheduling tool worth little to
a quiet practice may be valuable to a clinic losing appointments every day. A
repair that prevents one hour of downtime has different value on an idle
machine and on the constraint that stops an entire line. A premium material can
lower lifetime maintenance in one environment and be wasteful in another.
Price the consequence in context, then compare the offer with what that buyer
can really do instead.

A premium price can carry useful information when it supports materials,
reliability, service, curation, or a reputation earned through repeated
performance. Josh, the Austin e-commerce operator who compared favorable and
unfavorable reviews in Chapter 6, points to premium technology brands and urges
sellers to outvalue their alternatives. The defensible part of his method is
not copying a famous company's price. It is preserving what buyers repeatedly
value while repairing the failures they repeatedly experience. A high number
can signal trust after evidence exists; by itself, it can just as easily signal
scarcity theater, status, or an ambitious seller.

\section{A value story must earn belief}

At an Orlando event, a speaker asks what someone would pay for a small bag. He
offers it for \$1,000, then imagines that it contains \$100,000 in cash. The
arithmetic is easy: if the contents are certain, inspected, lawful, and truly
transfer with the purchase, the buyer receives \$99,000 more than the price.

But those conditions contain the whole problem. A stranger's claim that cash
may be hidden in a bag is not evidence. A sensible buyer asks what is inside,
then verifies the answer. Without inspection, the expected value depends on
the probability that the claim is true:

\begin{equation}
\mathbb{E}[V]=\sum_i p_i V_i.
\end{equation}

If the seller supplies no reason to assign a meaningful probability to the
valuable outcome, the hypothetical surplus cannot justify the purchase. The
speaker's insult toward people who resist the price is best discarded. Caution
is not a class defect. The useful habit is to look beyond price while remaining
equally curious about proof, alternatives, and downside.

Evidence changes value because it changes justified belief. A medical-device
entrepreneur named Jeremy describes years of conversations with doctors and
surgery centers, followed by three to six years of improving a neglected
product month by month. In his account, demand eventually began to arrive
without the same outbound effort. That history does not independently verify
the scale or medical outcomes he reports. It does illustrate a stronger basis
for an offer than a single forecast: repeated use, specific feedback,
successive improvement, and customers who seek the product.

Chris Meroff applies a harsher test to new ventures. When founders predict
millions of sales, he asks how many people have bought. He likes to expose an
idea to roughly a thousand conversations rather than rely on a persuasive
pitch. The number is a heuristic, not a universal sample requirement, and a
sale alone does not establish retention, margin, safety, or scalability. Still,
money paid under ordinary conditions is stronger evidence than enthusiasm
collected at no cost.

Build an evidence ladder for the particular outcome:

\begin{enumerate}[itemsep=0.15em]
\item \textbf{Mechanism.} Explain how the offer could cause the result and
  identify the step most likely to fail.
\item \textbf{Comparable history.} Show what happened in sufficiently similar
  settings, including failures and selection limits.
\item \textbf{Small proof.} Test the uncertain link at a safe scope with a
  baseline and an agreed measure.
\item \textbf{Use and recurrence.} Observe adoption, support burden, repeat
  use, renewal, or referral rather than counting the initial purchase alone.
\item \textbf{Attribution.} Separate the offer's contribution from changes in
  demand, staffing, season, policy, or other work performed at the same time.
\end{enumerate}

As evidence strengthens, the range of plausible value usually narrows. It may
narrow downward. That result is useful. A price reduced before a large rollout
can preserve a relationship; an offer whose value disappears under observation
should be redesigned or withdrawn.

\section{A price creates duties}

Pricing by value is sometimes taught as permission to extract as much as a
buyer can bear. That confuses willingness to pay with a complete ethical test.
A frightened patient, a desperate borrower, a bereaved family, or a company
facing a legal deadline may pay an enormous amount because refusal feels
impossible. Urgency can increase the economic value of competent help while
also increasing the seller's duty not to exploit impaired choice.

In Dubai, a medical-supply and promotional-products operator describes asking
questions until objections are exhausted, then making payment easy through
cards, credit, or invoicing. Questioning can uncover fit. It can also become a
way of steering a person past hesitation without resolving its cause. A clean
sale leaves the buyer able to say no. Payment convenience should remove
clerical friction, not hide total cost, financing charges, cancellation terms,
or the fact that the offer is unaffordable.

Another Dubai operator begins more constructively: understand the customer's
requirement, available money, and preference for a luxury or affordable
option, then explain benefits and return. Yet promised return must be supported
by the buyer's actual economics. A seller should not imply that an asset will
quickly double merely because such language closes a deal. Budget is also not
value. A buyer may have ample funds for a poor offer or insufficient cash for a
good one.

Four constraints belong beside every value claim:

\begin{itemize}[leftmargin=1.5em,itemsep=0.15em]
\item \textbf{Disclosure.} State material costs, uncertainties, conflicts,
  dependencies, and limits in language the buyer can use.
\item \textbf{Agency.} Preserve time, comparison, refusal, and informed consent;
  do not manufacture pressure that substitutes for evidence.
\item \textbf{Duty of care.} Never trade safety, legality, quality, privacy, or
  professional standards for a faster or more profitable outcome.
\item \textbf{Distribution.} Ask who receives the gain and who carries the
  burden. A payer's return can coexist with harm to users, workers, or the
  public.
\end{itemize}

These limits are not ornamental morality added after the calculation. They
change the calculation. Hidden implementation work, unsafe speed, transferred
risk, and damaged trust are costs even when the contract allows the seller to
ignore them.

\section{Write the value brief}

Before choosing a price, put the case on one page. The discipline is especially
useful when an attractive story is beginning to feel like proof.

\begin{enumerate}[itemsep=0.12em]
\item \textbf{Buyer and use.} Name the decision-maker, payer, users, affected
  people, setting, and period.
\item \textbf{Baseline.} Describe what happens without the offer, including the
  best available alternative and the option to wait.
\item \textbf{Outcome.} State the observable change, its unit, and when it should
  appear.
\item \textbf{Value range.} Estimate incremental gain, avoided loss, time and
  option value, then subtract implementation burden and downside risk.
\item \textbf{Evidence.} Record the mechanism, comparable results, proof already
  observed, uncertainty, and the fact that would most weaken the case.
\item \textbf{Delivery floor.} Include competent delivery, scarce capacity,
  support, expected rework, risk, and a sustainable margin.
\item \textbf{Price range.} Offer a range with assumptions rather than a precise
  number whose precision cannot be defended.
\item \textbf{Duties and review.} State disclosures, safeguards, conflicts,
  success measures, remedies, and the date at which evidence or price will be
  reconsidered.
\end{enumerate}

Then give the page to someone who can challenge it. Ask the buyer which term is
wrong. Ask the user which burden has been omitted. Ask the delivery team which
promise will be expensive to keep. Ask what a capable competitor would offer
instead. Revise the number when the evidence changes; a price is a decision,
not a confession of identity.

The value brief cannot close a sale. It can make the conversation honest enough
to deserve one. Price becomes real only when another person can question it,
refuse it, or buy.
